% ============================================================
% Problema 5
% ============================================================
\section{Problema 5: VLSM Inverso}

\begin{tcolorbox}[enunciado]
\textbf{Instrucciones:} Investiga qué es un \textbf{VLSM inverso} y resuelve el siguiente problema sin utilizar fórmulas. La solución debe obtenerse mediante el análisis visual de los bits de las direcciones IP.

\textbf{Enunciado:}\\
Un administrador de red encuentra cuatro subredes consecutivas ya configuradas en una sucursal:
\[
192.168.10.0/24, \quad 192.168.11.0/24, \quad 192.168.12.0/24, \quad 192.168.13.0/24
\]
Necesita resumirlas en una única ruta o \textbf{supernet} que abarque todas las subredes. Determina:
\begin{itemize}[leftmargin=*]
    \item El segmento de red (Network ID).
    \item La máscara en formato CIDR.
    \item La máscara en formato decimal.
    \item La primera IP de host utilizable.
    \item La última IP de host utilizable.
    \item La IP de broadcast.
    \item El número total de hosts utilizables.
\end{itemize}

\textbf{Restricción:} No utilices fórmulas para determinar la supernet. Representa el octeto que cambia en las cuatro direcciones mediante bits y encuentra visualmente qué bits permanecen constantes y cuáles cambian.

\textbf{Formato de respuesta:}
\vspace{0.3cm}
\begin{center}
    \begin{tabular}{lp{8cm}}
    \toprule
    \textbf{Parámetro} & \textbf{Resultado} \\
    \midrule
    Segmento de Red (ID) & \\
    Máscara CIDR & \\
    Máscara Decimal & \\
    Primer Host Útil & \\
    Último Host Útil & \\
    IP de Broadcast & \\
    Total de Hosts Utilizables & \\
    \bottomrule
    \end{tabular}
\end{center}
\vspace{0.3cm}

\textbf{Hint 1:} Observa únicamente el octeto que cambia entre las cuatro redes. Escríbelo en bits y alinea los cuatro valores verticalmente. Busca desde la izquierda los bits que son iguales en todas las direcciones. \\
\textbf{Hint 2:} Una vez encontrados los bits que permanecen constantes, imagina que todos los bits restantes son ``libres''. Ponlos en cero para encontrar el segmento de red y en uno para encontrar el broadcast. A partir de esos dos extremos puedes obtener el rango de hosts.
\end{tcolorbox}

\subsection*{Solución}
% Escribe aquí tu solución para el Problema 5...